\section{Versuchsaufbau und Durchführung}
\label{sec:aufbau}

\subsection{Aufbau}
Beschreibung des experimentellen Aufbaus. Eine Skizze gehört hierher:

\begin{figure}[H]
  \centering
  % \includegraphics[width=0.8\textwidth]{aufbau.pdf}
  \caption{Platzhalter: Foto/Skizze des Versuchsaufbaus.}
  \label{fig:aufbau}
\end{figure}

Verwendete Geräte und Materialien (Beispieltabelle):
\begin{table}[H]
  \centering
  \caption{Verwendete Geräte.}
  \label{tab:geraete}
  \begin{tabular}{lll}
    \toprule
    Gerät & Typ/Hersteller & Inventarnummer \\
    \midrule
    Oszilloskop & Mustertyp A & 12345 \\
    Multimeter & Mustertyp B   & 67890 \\
    Spannungsquelle & Mustertyp C & 11121 \\
    \bottomrule
  \end{tabular}
\end{table}

\subsection{Durchführung}
Schritt-für-Schritt-Beschreibung der Messung. Wichtig:
\begin{itemize}
  \item Alle Einstellungen (z. B. Verstärkung, Zeitbase) notieren.
  \item Systematische Fehler bereits hier erwähnen (z. B. Nullpunkt).
  \item Hinweis, wie oft und an welchen Stellen Messwiederholungen erfolgten.
\end{itemize}
